\section{When a Generated Answer Becomes a Result}

A fabricated reference can have a convincing title, plausible authors, and the right journal style. An invalid proof can have lemmas, notation, and impeccable transitions. A patch can look clean while calling an interface that does not exist. In each case, the output possesses a recognizable form. What it lacks is a completed route from that form to a dependable result.

Language models have made this separation industrially important. Candidate production can run ahead of candidate confirmation. The question is no longer simply where an answer came from. It is what has established the answer, at which level, and for which use.

\subsection{The Two Jobs of Form}

Form is a transferable constraint. A proof rule limits which steps may follow. A function signature limits which calls are well formed. A narrative structure limits what a reader expects next. Their constraints differ in strength and purpose, but each allows something learned in one setting to carry into another.

Formal logic isolates an especially powerful case. A valid inference preserves its validity when suitable contents are substituted into its structure. The achievement is compression: one rule can govern indefinitely many instances. The rule does not need to retain every instance separately.

The history from axiomatized geometry to symbolic logic made this power explicit. It also raised a harder ambition: could a sufficiently complete formal system settle every question in its domain and establish its own reliability from within?

The limits have precise statements. A consistent, effectively axiomatized theory strong enough for elementary arithmetic is incomplete. Under the usual conditions of the second incompleteness theorem, it cannot prove its own formalized consistency \cite{godel1931}. The halting problem rules out one algorithm deciding whether every arbitrary program halts \cite{turing1936}. Algorithmic information theory establishes that some finite strings have no substantially shorter description in a chosen universal description language \cite{chaitin1987}.

These results concern different objects. They do not show that every particular proof needs an unprovable assumption, that every program is beyond analysis, or that fluent language is inherently unreliable. What they prohibit are specific forms of universal closure. Engineers still prove particular properties of particular programs, using explicit assumptions and a specified logic.

Deployment adds a separate boundary: the correspondence between the formal object and the thing being built. A verified controller can satisfy a specification whose sensor model is wrong. A correct compiler can preserve the semantics of a program that expresses the wrong requirement. A sound proof checker can accept a theorem that is weaker than the informal claim used to advertise it.

The seL4 verification demonstrated a machine-checked refinement from an abstract specification to a C implementation, with explicit assumptions and exclusions \cite{klein2009sel4}. CompCert established semantic preservation for a substantial compiler, again within a defined account of source and target behavior \cite{leroy2009compcert}. These achievements make the boundary inspectable. They also show the work involved in connecting a claim to an implementation.

A result is valuable partly because one can say exactly what was established. “Verified” without the property, assumptions, and implementation boundary leaves the central question unanswered.

\subsection{How Prediction Learns Structure}

Earlier symbolic systems relied heavily on people writing rules. Learned models obtain much of their structure from records: language, code, diagrams, demonstrations, experimental reports, and the traces of other people's reasoning.

These records are highly patterned. Variable binding, references, function calls, nested expressions, and relations between goals and actions reduce uncertainty across many contexts. A model that represents such regularities can predict better than one that memorizes isolated fragments. Prediction can therefore create an incentive to learn useful abstract structure. Which structures it learns, and how reliably it uses them, remain empirical questions.

The physical process matters too. Training consumes electricity, moves data, changes stored states, and rejects heat. For an idealized reset of an unbiased bit, coupled to a thermal environment at temperature $T$, Landauer's bound is
\begin{equation}
 Q_{\min}=k_{\mathrm B}T\ln 2.
\end{equation}
At about $300\,\mathrm{K}$ this is roughly $2.9\times10^{-21}\,\mathrm{J}$ per bit reset. The bound concerns logically irreversible information erasure under the stated conditions; it is not a universal invoice attached to every arithmetic operation or memory write. Real hardware operates with additional device, communication, memory, and cooling costs.

Energy permits the computation. Structured data supplies learnable regularity. The objective and optimization process favor some representations over others. These roles cannot be collapsed into the claim that dissipation produces intelligence. A resistor dissipates energy without learning a grammar.

One can picture training as carving channels through a landscape: repeated data and error signals make some continuations easier than others. The image explains why widely recurring structures may become accessible together. It supplies no theorem that the resulting channels track truth.

The training objective rewards an operational target. In a prediction task, that target is predictive performance on a distribution of records. A true statement can be surprising in that distribution. A false statement can be conventional. A model can learn useful reasoning procedures and still apply them incorrectly to a particular case.

The GPT-4 technical report is a historical example of this coexistence: substantial capability accompanied by factual and reasoning failures \cite{openai2023gpt4}. It is evidence about a particular generation of systems, rather than a current frontier benchmark. The general engineering question persists as systems improve: what evidence connects this output to the object, rule, or environment on which it makes a claim?

Hallucinations have multiple sources, including data quality, objectives, retrieval failures, decoding, and reasoning errors \cite{huang2023hallucination}. Treating them as a single irreducible price of creativity would overstate the case. Individual error classes can be reduced, detected, or excluded by design. The harder problem is establishing reliability across the full range of uses advertised for a system.

\subsection{Three Places Where Confirmation Happens}

A bibliographic reference needs an identifiable source. A mathematical derivation needs valid inference. A warehouse plan needs an executable and acceptable intervention. These are three useful places to inspect a confirmation chain.

\begin{enumerate}
 \item \textbf{Reference.} Does the representation point to the claimed object? A paper title, API name, data field, or quoted finding must be traceable to something that supports its use.
 \item \textbf{Inference and execution.} Do the transformations obey the relevant rules? A proof, calculation, program, or causal argument needs the kind of check appropriate to its claim.
 \item \textbf{Action.} Does the proposal work in the receiving environment, with its resources, permissions, disturbances, and consequences?
\end{enumerate}

These checks can overlap and recur. They are not a mandatory chronology for every task. A program may run before anyone notices that it solves the wrong problem. An experiment may force revision of the representation used to design it.

AlphaGeometry offered an instructive architecture: a neural component proposed useful constructions and a symbolic engine performed geometric deduction \cite{trinh2024alphageometry}. Its published 2024 evaluation is a historical benchmark, not an assertion about the best available mathematics system in 2026. The interesting point is architectural: search and checking can have different mechanisms and different costs.

Software issue resolution offers a second judge. SWE-bench and its curated variants place a proposed change into a repository with an issue, dependencies, and tests \cite{swebench2024}. Passing the evaluation establishes something more specific than “can write code.” It still leaves questions about test coverage, security, maintainability, and production behavior.

Causal claims need another distinction. A relationship in observational data does not by itself specify the result of an intervention. Causal analysis makes assumptions about interventions and counterfactuals explicit \cite{pearl2009causality}. A persuasive causal story becomes actionable only when the assumptions and evidence support the proposed change.

The cost of traversing a chain can be expressed with conditional pass rates. If $p_i$ is the probability of passing stage $i$ \emph{conditional on having passed the preceding stages}, then
\begin{equation}
 p_{\mathrm{all}}=\prod_{i=1}^{k}p_i,
 \qquad
 N_{\mathrm{cand}}\approx
 \frac{M_{\mathrm{use}}}{\prod_{i=1}^{k}p_i}.
\end{equation}
The second expression is an expectation or planning approximation under stable, comparable trials. It does not guarantee a particular number of successes. Conditional pass rates require no independence assumption. If only unconditional rates are available, multiplying them requires additional justification.

The expression explains one possible migration of cost. Making proposals cheaper may leave experiments, trials, integration, and review as the expensive part of producing a usable result. Better models can improve those stages too: they can raise pass rates, design more informative tests, or reduce rework. Whether the total bill falls is a question about the entire chain.

Tools add confirmation interfaces. Retrieval can locate a source; an interpreter can execute a calculation; a proof assistant can check a formal derivation; a sensor can report a response. Each interface has a contract of its own. A search result may be stale. A test may miss the failure mode. A formally correct theorem may omit a necessary premise.

Reliability therefore belongs to the assembled system. A model that calls tools, reads results, revises its plan, and stops when evidence is inadequate has a different capability from the same model producing one unchecked response. The tools do not merely decorate its intelligence; they can supply part of the intelligence's working structure.

What remains costly varies by medium. A text can often be revised cheaply. A production migration needs monitoring and recovery. A medical intervention can require long follow-up, and an ecological intervention may have consequences that cannot be recalled. The verification ladder becomes heavier as the claim reaches deeper into the world.

Form makes knowledge portable. Confirmation determines where that portability can be trusted. A capable system needs both a way to generate candidates and a route by which a candidate becomes an established result.
